% =============================================================================
% БИЛЕТ 13
% =============================================================================

\ticket{13}{}

\question{Необходимое условие условного экстремума}

\begin{defn}[Точка условного локального экстремума]
  Точка $p$ называется \emph{точкой условного локального минимума} $f$ на $M$, если
  $\exists\delta>0$: $f(x) \geq f(p)$ для всех $x \in M \cap B_\delta(p)$.
\end{defn}

\begin{defn}[Функция Лагранжа]
  \emph{Функция Лагранжа} --- это $\mathcal{L}(x, \lambda) = f(x) - \lambda_1 p_1(x) - \ldots - \lambda_{n-k} p_{n-k}(x)$.

  На допустимом множестве $M = \{x: p_i(x) = 0\}$ выполняется
  $\mathcal{L}(x,\lambda) = f(x)$ для любого $\lambda$.
\end{defn}

\begin{theorem}[Необходимое условие условного экстремума, метод множителей Лагранжа]
  Пусть $p_1,\ldots,p_{n-k} \in C^1(\Omega)$, $M = \{x\in\Omega: p_1(x)=\ldots=p_{n-k}(x)=0\}$,
  $\forall p \in M$ градиенты $\nabla p_1(p),\ldots,\nabla p_{n-k}(p)$ линейно независимы
  и $f \in C^1(\Omega)$. Если $p$ --- точка условного экстремума $f$ на $M$, то
  $\exists \lambda_1^0,\ldots,\lambda_{n-k}^0$ такие, что
  \[
    \nabla f(p) + \sum_{j=1}^{n-k} \lambda_j^0\,\nabla p_j(p) = 0.
  \]
  \begin{proof}
    \textbf{Шаг 1.} Возьмём произвольный вектор $h \in T_p M$ (касательное пространство).
    По определению касательного вектора $\exists \gamma: (-\delta,\delta) \to M$
    такая, что $\gamma(0) = p$, $\gamma'(0) = h$.

    \textbf{Шаг 2.} Рассмотрим $g(t) = f(\gamma(t))$.
    Так как $p = \gamma(0)$ --- точка условного экстремума на $M$, то $t=0$ --- точка экстремума $g$,
    откуда $g'(0) = 0$. По правилу дифференцирования сложной функции:
    \[
      g'(0) = \langle \nabla f(\gamma(0)),\, \gamma'(0) \rangle = \langle \nabla f(p),\, h \rangle = 0.
    \]

    \textbf{Шаг 3.} Так как $h \in T_p M$ произвольно, $\nabla f(p) \perp T_p M$.

    \textbf{Шаг 4.} Поскольку ограничения задают многообразие $M$, каждый $\nabla p_j(p)$
    также перпендикулярен $T_p M$ (как нормаль к гиперповерхности $p_j = 0$).
    Так как $\{\nabla p_j(p)\}$ линейно независимы и порождают нормальное пространство
    (размерности $n-k$), то $\nabla f(p)$ лежит в их линейной оболочке:
    \[
      \nabla f(p) = -\sum_{j=1}^{n-k} \lambda_j^0\,\nabla p_j(p). \qedhere
    \]
  \end{proof}
\end{theorem}

\question{Линейность определённого интеграла. Аддитивность относительно отрезков интегрирования}

\begin{defn}[Сумма Римана]
  Пусть $f: [a,b] \to \mathbb{R}$, $T = \{x_0,\ldots,x_n\}$ --- разбиение, $\varepsilon_i \in [x_{i-1}, x_i]$.
  \emph{Сумма Римана}:
  $\sigma(f, T, \varepsilon_T) = \sum_{i=1}^n f(\varepsilon_i)(x_i - x_{i-1})$.
\end{defn}

\begin{defn}[Интегрируемость по Риману]
  $f \in R([a,b])$, если $\exists J \in \mathbb{R}$ такое, что
  $\forall\varepsilon>0\;\exists\delta>0: \forall T,\,\lambda(T)<\delta,\;\forall\varepsilon_T:
  |\sigma(f,T,\varepsilon_T) - J| < \varepsilon$.
  Обозначается $J = \int_a^b f(x)\,dx$.
\end{defn}

\begin{theorem}[Линейность]
  Если $f, g \in R([a,b])$ и $\alpha,\beta \in \mathbb{R}$, то $\alpha f + \beta g \in R([a,b])$ и
  \[
    \int_a^b (\alpha f + \beta g)\,dx = \alpha \int_a^b f\,dx + \beta \int_a^b g\,dx.
  \]
  \begin{proof}
    Из сумм Римана видно, что $\alpha$ и $\beta$ можно вынести за знак суммы:
    $\sigma(\alpha f + \beta g, T, \varepsilon_T) = \alpha\,\sigma(f,T,\varepsilon_T) + \beta\,\sigma(g,T,\varepsilon_T)$.
    Так как $\int f = J_1$ и $\int g = J_2$ по определению,
    для данного $\varepsilon > 0$ выбираем $\delta$ так, чтобы оба отклонения были $< \varepsilon / (|\alpha|+|\beta|)$,
    тогда $|\sigma(\alpha f+\beta g,T,\varepsilon_T) - \alpha J_1 - \beta J_2| < \varepsilon$. \qedhere
  \end{proof}
\end{theorem}

\begin{theorem}[Аддитивность]
  Если $a < c < b$, то
  \[
    \int_a^b f(x)\,dx = \int_a^c f(x)\,dx + \int_c^b f(x)\,dx.
  \]
  \begin{proof}
    Пусть $T_1$ --- разбиение $[a,c]$, $T_2$ --- разбиение $[c,b]$, $T = T_1 \cup T_2$.
    Тогда $\lambda(T_i) \leq \lambda(T)$ и
    $\sigma(f, T, \varepsilon_T) = \sigma(f, T_1, \varepsilon_{T_1}) + \sigma(f, T_2, \varepsilon_{T_2})$.
    При $\lambda(T) < \delta$ каждый из трёх интегралов имеет оценку отклонения $< \varepsilon$;
    из неравенства треугольника $|\sigma(f,T) - J_1 - J_2| < 3\varepsilon$. \qedhere
  \end{proof}
\end{theorem}
